\documentclass[11pt,a4paper]{article}
\usepackage[margin=1in]{geometry}
\usepackage{graphicx}
\usepackage{hyperref}
\usepackage{enumitem}
\usepackage{titlesec}
\usepackage{fancyhdr}

% Header and footer
\pagestyle{fancy}
\fancyhf{}
\rhead{Week 3 Report}
\lhead{MOT using Classical ML}
\rfoot{Page \thepage}

\title{\textbf{Weekly Progress Report - Week 3} \\
	\large Improve existing online trackers by measuring feature evolution of objects using
	Classical Machine Learning models [UAV Videos]}

\author{Group: Epochalypse\\\{Aagam Sheth, Mahima Parekh, Aakanksha Jadhav, Vansh Lilani\}}

\date{Reporting Period: March 2 - March 7, 2026 \\
	Submission Date: March 7, 2026}

\begin{document}
	
	\maketitle
	
	\section*{Outline of Performed Tasks:}
	\begin{itemize}[leftmargin=*]
		\item Implemented a modular framework to replace deep ReID embeddings with handcrafted classical features
		\item Analyzed tracker-specific implementations (ByteTrack, DeepSORT, BoT-SORT)
	\end{itemize}
	
	\section*{Dataset Update for MOT Experiments}
	The dataset used for the project has been changed to the \textbf{VisDrone dataset}. 
	This change was made because the previously considered \textbf{AU Drone dataset} primarily contains individual images rather than continuous video sequences.
	
	Since multi-object tracking (MOT) algorithms require sequential frames to track object identities over time, a dataset with video data is necessary. 
	The VisDrone dataset provides UAV-captured video sequences with annotations suitable for tracking tasks, making it more appropriate for evaluating MOT algorithms.
	
	\section*{BoxMOT Framework}
	We use the BoxMOT framework as the base system for implementing and testing different multi-object tracking algorithms. 
	BoxMOT provides a modular structure that connects object detectors (e.g., YOLO) with tracking algorithms and allows easy 
	integration of different appearance models. This flexibility makes it suitable for our work, where we aim to experiment 
	with replacing deep ReID embeddings with classical handcrafted features.
	
	\section*{Reference Materials Studied}
	
	\begin{enumerate}[leftmargin=*]		
		\item \textbf{ByteTrack}
		
		Studied the ByteTrack tracking method, which associates detections using motion prediction and considers both 
		high-confidence and low-confidence detections to improve tracking recall.
		
		\item \textbf{DeepSORT}
		
		Reviewed DeepSORT, which extends the SORT tracker by adding appearance-based ReID embeddings to improve identity 
		association during occlusions and long-term tracking.
		
		\item \textbf{BoT-SORT}

		Explored BoT-SORT, which improves upon ByteTrack by incorporating stronger ReID features and additional refinements 
		in the data association process.
	\end{enumerate}
	
	\section*{Tentative Tasks for Next Week}
	\begin{enumerate}
		\item Implement a baseline modular ReID feature extraction module (e.g., color histogram + HOG) 
		\item Integrate it into one tracker (ByteTrack or DeepSORT) to replace the deep ReID embedding
	\end{enumerate}
	
\end{document}